\section{Bisect Algorithmic Maps as a Reference Distribution}
\label{sec:behavior}

\subsection{The Reference Distribution Concept}

Post-\textit{Rucho} state court litigation requires courts to distinguish between
two sources of partisan skew in a redistricting plan: (1) skew produced by the
geographic sorting of partisan voters into urban-rural patterns---a structural
feature of American political geography that would appear in any map drawn
without partisan intent---and (2) skew produced by deliberate mapmaker choices
to pack or crack partisan communities. Federal and state doctrine uniformly holds
that the first source is permissible (mapmakers cannot control where voters
choose to live) while the second is the target of reform.

The evidentiary challenge is isolating the second source. Without a reference
distribution---a set of maps drawn by a neutral process from the same geographic
constraints---it is impossible to know how much partisan skew is ``natural'' for
a given state and how much reflects deliberate manipulation. Courts have struggled
with this challenge because the traditional counterfactual (``what would a fair
map look like?'') is contested: critics of any proposed neutral reference argue
that the reference itself embeds implicit partisan choices.

\textsc{Bisect} algorithmic maps resolve this challenge by providing a reference
distribution with a fully auditable, non-partisan generative process. The
\textsc{bisect} algorithm receives only census population counts and geographic
adjacency data; it cannot access voter registration files, election returns, or
any other partisan information. It is architecturally incapable of gerrymanders
because it lacks the inputs that gerrymandering requires. Any partisan properties
of \textsc{bisect} maps therefore reflect the underlying geographic distribution
of voters---the structural sorting component---rather than deliberate manipulation.

\subsection{How Bisect Maps Cluster on Partisan Metrics}

Under the \textsc{bisect} algorithm with geographic weights
(\texttt{--weights-override geographic}), maps produced for North Carolina
($k=14$), Wisconsin ($k=8$), Texas ($k=38$), and Florida ($k=28$) cluster
near the partisan-neutral centroid on all six L-series metrics:$^\dagger$

\begin{itemize}
\item \textbf{EG}: $|\text{EG}| < 0.03$ for NC ($k=14$) vs.\ enacted NC $|\text{EG}| = 0.09$.
  The threefold reduction is attributable to compactness-driven boundary-drawing
  that avoids the geographic segregation required for effective packing or cracking.
\item \textbf{MM}: bisect NC $\text{MM} \approx 0.02$ vs.\ enacted NC $\text{MM} \approx 0.07$.
  The residual $\text{MM} \approx 0.02$ reflects genuine geographic sorting of
  Democratic voters into the Triangle and Charlotte urban cores.
\item \textbf{Bias}: bisect NC $|\text{Bias}| < 0.02$ vs.\ enacted NC $\text{Bias} \approx -0.07$.
  Near-zero bias reflects the approximately symmetric treatment of Democratic
  and Republican voters under compactness optimization.
\item \textbf{Declination}: bisect NC $|\delta| < 0.10$ vs.\ enacted NC $\delta \approx +0.35$.
  The small residual in bisect maps reflects demographic concentration, not
  algorithmic cracking.
\item \textbf{Responsiveness}: bisect NC $R \approx 2.0$ vs.\ enacted NC $R \approx 1.3$.
  Compactness-driven bisection produces a roughly uniform distribution of
  district vote shares, maximizing competitive responsiveness.
\item \textbf{Gallagher LSq}: bisect NC $\text{LSq} \in [0.03, 0.08]$ vs.\
  enacted NC $\text{LSq} \in [0.10, 0.18]$. Single-member districts structurally
  preclude $\text{LSq} = 0$; the bisect floor reflects the geographic constraint.
\end{itemize}

\footnotetext{$^\dagger$ All single-run results. Values represent a single
\textsc{bisect} run per state using VEST 2020 presidential precinct returns
mapped to 2020 census tract geometry.}

\subsection{Why Compactness Produces Partisan Near-Neutrality}

The partisan near-neutrality of \textsc{bisect} maps is not a design goal---it
is a byproduct of compactness optimization applied to a geographically structured
electorate. The mechanism operates as follows.

Effective partisan gerrymandering requires creating districts with systematically
lopsided partisan compositions: packing the opposing party's voters into a few
districts with 80\%+ supermajorities (wasting their excess votes) and cracking
their remaining voters across many districts at 35-45\% (wasting their losing
votes). Both operations require drawing geographically elongated or irregular
boundaries that reach across natural geographic discontinuities to aggregate
partisans separated by space.

Compactness optimization penalizes exactly the boundary configurations that
packing and cracking require. A compact district is, by definition, one whose
boundary does not extend into distant territory to collect partisan allies;
its composition is determined primarily by who lives nearby. Because American
residential geography has a fractal property---partisan communities are mixed
at finer scales than they appear at coarser scales---compactness constraints
tend to produce districts that are competitive near the geographic center of
the state's partisan distribution.

This mechanism explains a key empirical finding that appears consistently across
NC, WI, TX, and FL: \textsc{bisect} maps reduce partisan skew on all five
descriptive metrics simultaneously, without any trade-off. A mapmaker attempting
to reduce EG at the expense of Bias, or to improve MM at the expense of
Declination, would need to engage in the very geographic reaching that
compactness prohibits. Compactness constraint therefore functions as a partisan
fairness constraint, not as an alternative to it.
